\chapter{Introduction}
\label{ch:intro}

Formal specifications buy precision that informal prose cannot match.
That precision becomes a liability when large language models (LLMs) write the
implementation: the model has no native notion of guard order, so it produces
code that is statistically plausible yet formally wrong.
The resulting faults pass most everyday inputs and break exactly where the
specification draws its boundaries.
Test-based repair loops miss them because they sample the input space without
knowing which regions matter.

The question that drives this report is therefore practical:
\emph{can specification-conformance defects in LLM-generated code be prevented
by using the specification itself as the primary evidence source, without full
deductive verification?}
The answer developed below is affirmative.
The vehicle is the \textbf{Specification-guided Defect Prevention Framework}
(\textsc{SgDP}), with \textbf{HSP-Agile} as its SOFL/FSF instantiation.


\section{Problem Statement}
\label{sec:intro:problem}

\subsection{Setting}
SOFL \citep{liu1997sofl,liu2014sofl} provides the Functional Specification
Format (FSF): an ordered list of scenarios, each pairing a guard with a
postcondition under first-match precedence.
Agile-SOFL \citep{li2022agilesofl} embeds those artefacts in sprint workflows
as design documents, oracles, and acceptance criteria.
LLM-assisted generation accelerates that workflow
\cite{chen2021codex,openai2023gpt4}, but synthesises by pattern completion
rather than by respecting guard order
\cite{dakhel2023copilot,lynn2024verifierloop}.
Overlapping guards therefore expose precedence-sensitive faults that standard
test suites miss \citep{zhu1997specification}.

\subsection{What Goes Wrong}
A specification-conformance defect is, in plain terms, a bug that unit tests
rarely catch.
The generated function may look correct on the happy path, yet on a thin slice
of the input space---often where two guards overlap---it returns the output of
the wrong scenario.
Random or distribution-biased tests almost never land in that slice.
A repair loop that only sees failing unit tests never learns about the ordering
constraint.

That observation leads to a release decision.
Given task~$t$ with ordered scenarios $\mathcal{S}_t = (S_1,\ldots,S_n)$ and
synthesis budget~$K$, the question is whether any candidate~$c$ satisfies
\begin{equation}
  \mathrm{Accept}(c, t)
  \;\iff\;
  \mathrm{Conf}(c, t) \geq \tau
  \;\wedge\;
  P^H(c) = 0,
  \label{eq:acceptance}
\end{equation}
where $\mathrm{Conf}(c,t)\in[0,1]$ is the fraction of specification-derived
witnesses on which $c$ matches the oracle, $\tau=1.0$ in this work, and
$P^H(c)=0$ asserts that no high-severity structural pattern remains.
A witness for scenario~$S_i$ must satisfy
$G_i(\mathbf{x}) \wedge \bigwedge_{j<i}\neg G_j(\mathbf{x})$---a conjunction
that standard sampling heuristics cannot reliably reach
\cite{barrett2018smt}.

\subsection{A Concrete Example}
The running example throughout the report is \code{classify\_signal}
(Listing~\ref{lst:fsf-spec}).
It is small enough to hold in working memory, yet rich enough to expose the
ordering conflict that motivates the whole framework.

\inputfsf{listings/classify_signal.fsf}{FSF specification for \code{classify\_signal} (pseudocode notation).}{lst:fsf-spec}

When \code{threshold} is negative, the input
$(\mathit{level}{=}{-}3,\;\mathit{threshold}{=}{-}10)$ satisfies both $G_1$
(\texttt{level < 0}) and $G_2$ (\texttt{level >= threshold}).
FSF ordering mandates \code{"Error"} via~$S_1$.
A typical LLM candidate evaluates the threshold comparison first:

\inputpython{listings/classify_signal_bug.py}{LLM-generated candidate with a scenario-ordering violation.}{lst:llm-incorrect}

The candidate returns \code{"Critical"} instead of \code{"Error"}.
A test suite that draws \code{threshold} from a non-negative distribution never
reaches this region---exactly the latent defect that
Equation~\eqref{eq:acceptance} must prevent.
Figure~\ref{fig:defect-vs-test} makes the sampling contrast concrete: denser
random tests do not help; a specification-aware witness does.

\begin{figure}[t]
  \centering
  \tikzinput{diagrams/tikz/defect_vs_test}
  \caption{\textbf{Why unit tests miss the defect.}
    Left: random tests concentrate on $\mathit{threshold}\ge 0$ and never
    exercise the ordering conflict.
    Right: an SMT witness for $\Phi_1$ forces a point in the missed region
    (e.g.\ $(-3,-10)$), exposing the conformance failure.
    The figure motivates specification-aware witnesses rather than denser
    random testing.}
  \label{fig:defect-vs-test}
\end{figure}


\section{Research Challenges}
\label{sec:intro:challenges}

Three obstacles stand between the example above and a usable prevention
pipeline.

\subsection{Closing the Specification--Code Semantic Gap}
Checking scenario~$S_i$ requires inputs in
$G_i \wedge \bigwedge_{j<i}\neg G_j$, often a low-measure region missed by
random sampling.
SMT solvers \citep{demoura2008z3,barrett2018smt} can generate such witnesses,
but only if the FSF encoding is faithful.
The response in \textsc{SgDP} is a constraint-guided witness generator that
encodes guard precedence as SMT constraints (Chapter~\ref{ch:method}).

\subsection{Structured Counterexample-Guided Repair}
Test-feedback loops lack specification-level structure
\cite{lynn2024verifierloop}: the model may patch one counterexample without
fixing guard precedence.
Effective repair needs prompts that name the violated scenario, state the
ordering constraint, and include the SMT witness.
\textsc{SgDP} therefore structures each failure into a scenario-aware record
\emph{before} rendering natural-language prompts---the Semantic Feedback IR
introduced in Chapter~\ref{ch:method}.

\subsection{Operationalising Prevention within a Bounded Budget}
Full verification is impractical for unannotated LLM output
\cite{leino2010dafny,loughman2024dafnybench}.
What is needed is calibrated prevention: stronger than test-feedback
\cite{chen2021codex}, weaker than deductive proof, and fast enough for sprint
budgets \citep{beck2001xp,li2022agilesofl}.
\textsc{SgDP} answers with a conjunctive acceptance predicate and a hard repair
budget~$K$ ($K=3$ in the experiments).


\section{Proposed Approach}
\label{sec:intro:approach}

HSP-Agile runs a four-stage loop on each task.
It synthesises a candidate from the FSF text, checks that candidate against SMT
witnesses derived from the specification, repairs with scenario-aware feedback
when checks fail, and gates release with a pattern screen.
Figure~\ref{fig:sgdp-overview} is the map of that loop: adapters feed SpecIR,
witnesses and Semantic Feedback IR drive repair, and acceptance is conjunctive.
Chapter~\ref{ch:method} develops the general \textsc{SgDP} interface, states
three correctness results, and walks through \code{classify\_signal} end to
end; Chapter~\ref{ch:results} shows that the same loop transfers to Mini-Z and
Mini-StateMachine via a SpecIR adapter layer.

\begin{figure}[t]
  \centering
  \tikzinput{diagrams/tikz/sgdp_framework_overview}
  \caption{\textbf{\textsc{SgDP} framework overview.}
    Surface specification text is parsed by a notation adapter into canonical
    SpecIR; SMT witnesses and a JSON-serialised Semantic Feedback IR drive
    counterexample-guided LLM repair under budget~$K$; hybrid acceptance gates
    release on formal conformance and pattern safety.
    The dashed branch shows the HSP-Agile FSF lowerer and pattern guard.
    Read left to right for one attempt; the dashed orange arc is the repair
    loop.}
  \label{fig:sgdp-overview}
\end{figure}


\section{Contributions}
\label{sec:intro:contributions}

Five contributions structure the rest of the report
(Figure~\ref{fig:contribution-map}).

\begin{figure}[t]
  \centering
  \tikzinput{diagrams/tikz/contribution_map}
  \caption{Overview of \textsc{SgDP} research contributions.
    The five contributions form an integrated path from pipeline design through
    empirical validation.}
  \label{fig:contribution-map}
\end{figure}

\begin{enumerate}
  \item[\textbf{C1.}] \textbf{Specification-guided Defect Prevention Framework
    (\textsc{SgDP}).}
    A notation-agnostic SpecIR layer with pluggable adapters and a four-stage
    prevention loop; HSP-Agile is the SOFL/FSF instantiation.

  \item[\textbf{C2.}] \textbf{Strict Acceptance Predicate and Semantic Feedback
    IR.}
    The release predicate $\mathrm{Accept}(c,t)$ and a nine-field JSON record
    that separates failure structure from prompt rendering.

  \item[\textbf{C3.}] \textbf{Precedence-Sensitive Hard Benchmark.}
    120 synthetic tasks with overlapping guards, plus real-derived and
    multi-notation subsets.

  \item[\textbf{C4.}] \textbf{Two-Layer Prevention Evaluation Protocol.}
    PDR/FAR measurement via specification-confusion and implementation-screening
    mutants.

  \item[\textbf{C5.}] \textbf{Empirical Evidence and Reproducible Artifacts.}
    90.4\% mean strict conformance (+6.2\,pp over one-shot LLM); full logs and
    scripts released.
\end{enumerate}


\section{Research Questions}
\label{sec:intro:rq}

Five research questions organise the evaluation.
Results appear in \emph{mechanism-first} order in Chapter~\ref{ch:results}
(why the framework works, then that it works overall); the numbering below
follows the experimental logic of Chapter~\ref{ch:experiments}.

\textbf{RQ1}---overall effectiveness of specification-guided prevention (E1, E2);
\textbf{RQ2}---fault-class coverage and residual failure modes (E2, E7, E9);
\textbf{RQ3}---component attribution via ablation and feedback variants (E1, E5, E6);
\textbf{RQ4}---scaling across complexity dimensions (E3, E4);
\textbf{RQ5}---quality--overhead trade-off (E1).


\section{Report Organisation}
\label{sec:intro:organisation}

The remainder of the report follows the path from background to evidence.
Background material on Agile-SOFL, LLM synthesis, SMT, and mutation testing
appears in Chapter~\ref{ch:background}.
Chapter~\ref{ch:formalization} makes the conformance problem precise on
\code{classify\_signal}; Chapter~\ref{ch:method} turns that problem into the
\textsc{SgDP} method with a worked end-to-end example and formal theorems.
The software realisation and experimental protocol occupy
Chapters~\ref{ch:implementation} and~\ref{ch:experiments}.
Results across all five research questions---including generalisation to
Mini-Z and Mini-StateMachine (Section~\ref{ch:generalization})---appear in
Chapter~\ref{ch:results}; interpretation and threats follow in
Chapter~\ref{ch:discussion}.
Chapter~\ref{ch:conclusion} closes with contributions and future directions.
Three appendices cover reproducibility
(Appendix~\ref{app:reproducibility}), benchmark construction
(Appendix~\ref{app:benchmark}), and the fault pattern catalogue
(Appendix~\ref{app:patterns}).
